% =====================================================================
\chapter{Resultados e discussão}\label{cap:resultados}
% =====================================================================
% Todos os números deste capítulo saem de tcc/analise/resultados_c3.py, que lê a base
% da Economática e gera as tabelas em tabelas/*.tex. Se a base mudar, rode o script de
% novo e confira os números citados no texto (procure por "%" seguido de "fonte:").

Este capítulo apresenta a evolução do endividamento dos dois grupos, os testes de
diferença no período completo e nos subperíodos e as análises de robustez, e
discute os resultados à luz do referencial teórico.

% ---------------------------------------------------------------------
\section{Evolução do endividamento dos grupos}\label{sec:evolucao}
% ---------------------------------------------------------------------

A \autoref{tab:anuais} apresenta as médias anuais dos três indicadores centrais da
análise. No endividamento oneroso, medido pela Dívida bruta / Ativo total, as
empresas familiares ficaram abaixo das não familiares em todos os oito exercícios.
A distância, porém, não foi constante: em 2018 e 2019, os grupos estavam próximos
(23,9\% contra 26,4\% e 27,0\% contra 28,3\%); a diferença se ampliou em 2021 e 2022,
quando as não familiares chegaram a 35,4\% e 36,1\%, contra 26,0\% e 24,9\% das
familiares; e voltou a se estreitar a partir de 2024, até quase desaparecer em 2025
(26,4\% contra 26,9\%).
% fonte: saidas/resultados.md, seção 2

% Gerado por tcc/analise/resultados_c3.py -- não edite à mão; rode o script de novo.
\begin{table}[htb]
\caption{Médias anuais por grupo em 31 de dezembro (\%)}\label{tab:anuais}
\centering\footnotesize\setlength{\tabcolsep}{4pt}
\begin{tabular}{lrrrrrr}
\toprule
Ano & \multicolumn{2}{c}{Dívida bruta / Ativo} & \multicolumn{2}{c}{Exigível / Ativo} & \multicolumn{2}{c}{Dívida CP / Dívida total} \\
\midrule
 & F & NF & F & NF & F & NF \\
2018 & 23,9 & 26,4 & 50,2 & 47,8 & 29,2 & 41,4 \\
2019 & 27,0 & 28,3 & 56,1 & 53,3 & 23,0 & 25,4 \\
2020 & 29,6 & 30,5 & 59,0 & 56,0 & 17,2 & 25,8 \\
2021 & 26,0 & 35,4 & 54,5 & 60,1 & 21,9 & 47,0 \\
2022 & 24,9 & 36,1 & 54,9 & 59,0 & 19,6 & 35,9 \\
2023 & 26,2 & 32,8 & 55,5 & 56,8 & 18,9 & 36,8 \\
2024 & 24,7 & 28,2 & 55,6 & 54,6 & 16,1 & 39,1 \\
2025 & 26,4 & 26,9 & 57,2 & 62,1 & 19,1 & 54,3 \\
\bottomrule
\end{tabular}
\fonte{Elaborado pelo autor (2026), com dados da Economática.}
\nota{F: familiares (10 empresas); NF: não familiares (6 empresas; 5 na Dívida CP / Dívida total). Média simples entre as empresas de cada grupo.}
\end{table}


O endividamento geral, medido pelo Exigível total / Ativo total, não seguiu o
mesmo padrão. As familiares tiveram médias superiores às das não familiares em
quatro dos oito exercícios, incluindo os anos anteriores à pandemia, e as médias
dos dois grupos ficaram entre 47,8\% e 62,1\% ao longo de todo o período. A
diferença observada na dívida onerosa, portanto, não corresponde a um passivo
total menor, e sim a uma composição do passivo com menos dívida financeira.

A composição da dívida por prazo apresentou a diferença mais persistente. Em todos
os exercícios, a participação da dívida de curto prazo foi menor nas familiares, e
a distância se tornou expressiva a partir de 2021: desde então, a dívida de curto
prazo representou entre 16\% e 22\% da dívida total das familiares, contra 36\% a
54\% nas não familiares.

% Gráficos do autor (exportar do Google Planilhas para figuras/ como PDF ou PNG):
% \begin{figure}[htb]
%   \caption{Dívida bruta / Ativo total: média trimestral por grupo (2018--2025)}\label{fig:divbruta}
%   \centering
%   \includegraphics[width=\textwidth]{figuras/div-bruta-ativo.pdf}
%   \fonte{Elaborado pelo autor (2026), com dados da Economática.}
% \end{figure}
% (repetir para Exigível / Ativo e Dívida CP / Dívida total)
\verificar{inserir os gráficos trimestrais dos três indicadores (Figuras 1 a 3) e
citá-los neste ponto do texto}

% ---------------------------------------------------------------------
\section{Comparação entre os grupos no período completo}\label{sec:completo}
% ---------------------------------------------------------------------

A \autoref{tab:completo} apresenta os testes de diferença para o período de 2018 a
2025. Nos cinco indicadores, a média das familiares foi inferior à das não
familiares, mas nenhuma diferença foi significativa ao nível de 5\% em qualquer dos
dois testes.

% Gerado por tcc/analise/resultados_c3.py -- não edite à mão; rode o script de novo.
\begin{table}[htb]
\caption{Endividamento das empresas familiares (F) e não familiares (NF), 2018--2025}\label{tab:completo}
\centering\footnotesize\setlength{\tabcolsep}{4pt}
\begin{tabular}{lrrcccc}
\toprule
Indicador & F & NF & n & Anos F $<$ NF & p (t) & p (MW) \\
\midrule
Dívida bruta / Ativo total & 26,1\% & 30,6\% & 10 $\times$ 6 & 8 de 8 & 0,556 & 0,255 \\
Exigível total / Ativo total & 55,4\% & 56,2\% & 10 $\times$ 6 & 4 de 8 & 0,885 & 0,957 \\
Dívida bruta / (Dívida bruta + PL) & 37,2\% & 43,5\% & 10 $\times$ 6 & 6 de 8 & 0,563 & 0,255 \\
Dívida líquida / EBITDA & 1,4x & 4,4x & 10 $\times$ 6 & 5 de 8 & 0,285 & 0,416 \\
Dívida CP / Dívida total & 20,4\% & 38,2\% & 10 $\times$ 5 & 8 de 8 & 0,112 & 0,058$^{*}$ \\
\bottomrule
\end{tabular}
\fonte{Elaborado pelo autor (2026), com dados da Economática.}
\nota{Testes bicaudais sobre a média de cada empresa nas datas de 31/12 do período. p (MW): teste de Mann-Whitney, aproximação normal com correção de continuidade. $^{***}$, $^{**}$ e $^{*}$: significativo a 1\%, 5\% e 10\%. Valores: média simples das médias por empresa. Na Dívida CP / Dívida total, a Bradsaúde é excluída por não ter dívida financeira; na Dívida líquida / EBITDA, a Biomm tem um único ano com EBITDA positivo (2025).}
\end{table}


No indicador principal, a Dívida bruta / Ativo total, a média das familiares foi
de 26,1\%, contra 30,6\% das não familiares, diferença de 4,5 pontos percentuais que
não foi estatisticamente significativa (p = 0,556 no teste t e p = 0,255 no teste
de Mann-Whitney). O Exigível total / Ativo total praticamente não diferiu entre os
grupos (55,4\% e 56,2\%; p = 0,885 e 0,957). A estrutura de capital, medida pela
Dívida bruta / (Dívida bruta + PL), seguiu a mesma direção da dívida onerosa
(37,2\% contra 43,5\%), também sem significância.

O resultado mais próximo da significância foi o da composição da dívida: a dívida
de curto prazo representou, em média, 20,4\% da dívida total das familiares e
38,2\% da dívida das não familiares. A diferença foi significativa apenas ao nível
de 10\% no teste de Mann-Whitney (p = 0,058) e não foi significativa no teste t
(p = 0,112). A divergência entre os testes é coerente com a dispersão do grupo não
familiar, em que a Lifemed concentrou 66\% de sua dívida no curto prazo, valor que
eleva a variância e reduz o poder do teste t.
\verificar{conferir no Google Planilhas o valor da LMED (média de 66,3\%) antes de
citar}

Na Dívida líquida / EBITDA, a média das não familiares (4,4x) foi três vezes a das
familiares (1,4x), mas esse contraste se deve, em grande parte, à Biomm, cujo
indicador só pôde ser calculado em 2025 (16,1x), já que o EBITDA da empresa foi
negativo nos anos anteriores. Sem a Biomm, a média das não familiares cai para
2,0x (p = 0,560 no teste t), e as medianas dos grupos são 1,6x e 2,6x. Esse
indicador, portanto, não permite concluir que as familiares tenham maior
capacidade de pagamento.
\decidir{confirmar este tratamento da Biomm (apresentar o valor com e sem a empresa
e as medianas) e ajustar a seção 3.4 de acordo}

% ---------------------------------------------------------------------
\section{Subperíodos: antes, durante e depois da pandemia}\label{sec:subperiodos}
% ---------------------------------------------------------------------

A \autoref{tab:subperiodos} repete os testes em cada subperíodo. Antes da pandemia
(2018--2019), os grupos foram semelhantes em todos os indicadores, e nenhuma
diferença se aproximou da significância.

% Gerado por tcc/analise/resultados_c3.py -- não edite à mão; rode o script de novo.
\begin{table}[htb]
\caption{Comparação entre os grupos por subperíodo}\label{tab:subperiodos}
\centering\footnotesize\setlength{\tabcolsep}{4pt}
\begin{tabular}{llrrcc}
\toprule
Indicador & Subperíodo & F & NF & p (t) & p (MW) \\
\midrule
Dív. bruta / Ativo & 2018--2019 & 25,4\% & 27,3\% & 0,815 & 0,625 \\
 & 2020--2021 & 27,8\% & 33,0\% & 0,543 & 0,255 \\
 & 2022--2025 & 25,5\% & 31,0\% & 0,497 & 0,416 \\
Exigível / Ativo & 2018--2019 & 53,1\% & 50,6\% & 0,703 & 0,704 \\
 & 2020--2021 & 56,8\% & 58,1\% & 0,825 & 0,871 \\
 & 2022--2025 & 55,8\% & 58,1\% & 0,760 & 0,957 \\
Dív. CP / Dív. total & 2018--2019 & 24,9\% & 33,4\% & 0,435 & 0,298 \\
 & 2020--2021 & 19,5\% & 36,4\% & \textbf{0,040}$^{**}$ & \textbf{0,043}$^{**}$ \\
 & 2022--2025 & 18,4\% & 41,5\% & 0,130 & 0,098$^{*}$ \\
\bottomrule
\end{tabular}
\fonte{Elaborado pelo autor (2026), com dados da Economática.}
\nota{Testes bicaudais sobre a média de cada empresa nas datas de 31/12 do período. p (MW): teste de Mann-Whitney, aproximação normal com correção de continuidade. $^{***}$, $^{**}$ e $^{*}$: significativo a 1\%, 5\% e 10\%. Subperíodos: pré-pandemia (2018--2019), pandemia (2020--2021) e pós-pandemia (2022--2025).}
\end{table}


Durante a pandemia (2020--2021), a diferença na composição da dívida tornou-se
significativa ao nível de 5\% nos dois testes: a dívida de curto prazo representou
19,5\% da dívida total das familiares e 36,4\% da das não familiares (p = 0,040 no
teste t e p = 0,043 no teste de Mann-Whitney). No mesmo período, a dívida onerosa
das familiares também foi menor (27,8\% contra 33,0\% do ativo), mas sem
significância (p = 0,543 e 0,255), e o endividamento geral permaneceu semelhante
(56,8\% e 58,1\%).

Após a pandemia (2022--2025), a distância na composição da dívida aumentou em
termos absolutos (18,4\% contra 41,5\%), mas a significância se reduziu: a
diferença foi significativa apenas a 10\% no teste de Mann-Whitney (p = 0,098) e
não foi significativa no teste t (p = 0,130).

Como explicado na \autoref{sec:procedimentos}, a divisão em subperíodos tem caráter
exploratório, e o resultado da pandemia é o único significativo a 5\% entre os
testes realizados. Ele deve, por isso, ser lido como um indício compatível com a
hipótese, e não como sua confirmação.

% ---------------------------------------------------------------------
\section{Análises de robustez}\label{sec:robustez}
% ---------------------------------------------------------------------

A \autoref{tab:robustez} mostra como os resultados se alteram com a exclusão da
Bradsaúde e com as classificações alternativas de empresa familiar.

% Gerado por tcc/analise/resultados_c3.py -- não edite à mão; rode o script de novo.
\begin{table}[htb]
\caption{Análises de robustez}\label{tab:robustez}
\centering\footnotesize\setlength{\tabcolsep}{4pt}
\begin{tabular}{llrrccc}
\toprule
Especificação & Indicador & F & NF & n & p (t) & p (MW) \\
\midrule
C3 (principal) & Dív. bruta / Ativo & 26,1\% & 30,6\% & 10 $\times$ 6 & 0,556 & 0,255 \\
 & Exigível / Ativo & 55,4\% & 56,2\% & 10 $\times$ 6 & 0,885 & 0,957 \\
 & Dív. CP / Dív. total & 20,4\% & 38,2\% & 10 $\times$ 5 & 0,112 & 0,058$^{*}$ \\
C3 sem Bradsaúde & Dív. bruta / Ativo & 26,1\% & 36,7\% & 10 $\times$ 5 & \textbf{0,027}$^{**}$ & 0,058$^{*}$ \\
 & Exigível / Ativo & 55,4\% & 59,8\% & 10 $\times$ 5 & 0,348 & 0,759 \\
 & Dív. CP / Dív. total & 20,4\% & 38,2\% & 10 $\times$ 5 & 0,112 & 0,058$^{*}$ \\
C2 & Dív. bruta / Ativo & 28,7\% & 27,0\% & 7 $\times$ 9 & 0,801 & 1,000 \\
 & Exigível / Ativo & 53,6\% & 57,3\% & 7 $\times$ 9 & 0,567 & 0,597 \\
 & Dív. CP / Dív. total & 19,1\% & 32,6\% & 7 $\times$ 8 & 0,080$^{*}$ & 0,118 \\
C1 & Dív. bruta / Ativo & 30,9\% & 26,3\% & 5 $\times$ 11 & 0,558 & 0,734 \\
 & Exigível / Ativo & 55,8\% & 55,7\% & 5 $\times$ 11 & 0,987 & 0,910 \\
 & Dív. CP / Dív. total & 19,2\% & 29,9\% & 5 $\times$ 10 & 0,131 & 0,245 \\
\bottomrule
\end{tabular}
\fonte{Elaborado pelo autor (2026), com dados da Economática.}
\nota{Testes bicaudais sobre a média de cada empresa nas datas de 31/12 do período. p (MW): teste de Mann-Whitney, aproximação normal com correção de continuidade. $^{***}$, $^{**}$ e $^{*}$: significativo a 1\%, 5\% e 10\%. C1, C2 e C3: Cenários 1, 2 e 3; C1 e C2 são classificações alternativas de empresa familiar (\autoref{qua:cenarios}).}
\end{table}


A exclusão da Bradsaúde, única empresa da amostra sem dívida financeira, eleva a
média de endividamento oneroso das não familiares de 30,6\% para 36,7\% e torna a
diferença significativa a 5\% no teste t (p = 0,027) e a 10\% no teste de
Mann-Whitney (p = 0,058). O resultado do indicador principal depende, portanto, de
uma única observação: com ela, a diferença não é significativa; sem ela, passa a
ser.

As classificações alternativas produzem um quadro distinto. Nos Cenários 1 e 2, que
exigem participação familiar maior, a diferença na dívida onerosa inverte o sinal:
as familiares passam a ter média superior à das não familiares (30,9\% contra
26,4\% no Cenário 1 e 28,7\% contra 27,0\% no Cenário 2), sem significância. A
inversão ocorre porque RaiaDrogasil, Profarma e Hypera, familiares apenas no
Cenário 3, passam ao grupo não familiar. As duas primeiras, que atuam no comércio e
na distribuição de medicamentos, estão entre as empresas menos endividadas da
amostra (12,8\% e 16,5\% do ativo, em média). Parte da diferença observada no
Cenário 3 pode refletir, portanto, o modelo de negócio desse subsetor, e não
apenas a estrutura de controle.

A composição da dívida, por outro lado, manteve a direção nos três cenários: a
participação da dívida de curto prazo foi menor nas familiares em todos eles (19,1\%
contra 32,6\% no Cenário 2, com p = 0,080 no teste t, e 19,2\% contra 29,9\% no
Cenário 1, sem significância). É, entre os indicadores analisados, o resultado
menos sensível ao critério de classificação.

% ---------------------------------------------------------------------
\section{Discussão}\label{sec:discussao}
% ---------------------------------------------------------------------

Os resultados não permitem confirmar a hipótese H1. Embora as empresas familiares
tenham apresentado menor endividamento oneroso em todos os exercícios, a diferença
não foi estatisticamente significativa na amostra principal, não se manteve sob as
classificações alternativas e dependeu da presença de uma única empresa sem dívida.
Não foram encontradas, portanto, evidências estatisticamente significativas de que
as empresas familiares do setor de saúde da B3 se endividem menos que as não
familiares.

Esse resultado é coerente com a literatura inconclusiva apresentada no
\autoref{cap:referencial}. As duas forças discutidas na \autoref{sec:risco} podem
atuar em sentidos opostos: a aversão à interferência de credores, associada à
preservação da riqueza socioemocional \cite{santos2018,ames2020}, levaria a menos
dívida, enquanto o receio de diluir o controle levaria a preferir a dívida à emissão
de ações \cite{poletti2022}. No Brasil, \citeonline{kayo2018} encontraram empresas
familiares mais endividadas, o que indica que a segunda força pode ser relevante no
mercado brasileiro. Em uma amostra pequena, a combinação dessas forças é compatível
com diferenças de pequena magnitude e instáveis entre especificações, como as
observadas.

O resultado mais consistente refere-se ao prazo, e não ao nível, da dívida. As
familiares concentraram menos dívida no curto prazo em todos os exercícios e nos
três critérios de classificação, e a diferença foi significativa durante a
pandemia. Esse padrão é compatível com o argumento de \citeonline{diamond1991}: a
dívida de curto prazo precisa ser renovada com frequência e expõe a empresa ao
risco de que os credores não a refinanciem, risco que se agrava em períodos de
crise. Para uma família que busca preservar a autonomia decisória, alongar a dívida
reduz a necessidade de renegociações frequentes e, com ela, a exposição à avaliação
dos credores. Assim, a aversão à interferência de credores parece se manifestar
mais na escolha do prazo do que na quantidade de dívida.

Dois cuidados limitam essa interpretação. Primeiro, a diferença de prazo foi
significativa a 5\% apenas em um subperíodo definido durante a análise, entre
diversos testes realizados, o que aumenta a chance de um resultado significativo
ocorrer por acaso. Segundo, o desenho descritivo não permite separar o efeito do
controle familiar de outras características das empresas, como o subsetor de
atuação \cite{oliveira2020} e o tamanho. Os resultados devem, por isso, ser lidos como evidência
descritiva de uma postura mais conservadora das famílias em relação ao prazo da
dívida, a ser testada em amostras maiores e com controles.
